\begin{textsample}{POS Dim 1 – vem\_ed\_17 – Score 20.00 – vem\_ed\_17/t075.txt}  \label{ex:f1_pos_131}
Fazer COIL \textbf{permite} enxergar \textbf{que} há \textbf{pessoas} no mundo \textbf{que} \textbf{são} como eu, \textbf{permite} ver \textbf{questões} como as consequências da escravidão no Brasil e nos Estados Unidos.
Os estudantes merecem trazer sua cultura para a sala de aula. \textbf{Parte} do \textbf{que} vamos fazer nos próximos cinco anos envolve o pensamento inclusivo e a interseccionalidade.
Como falamos sobre \textbf{ações} sociais, \textbf{ambientais} e econômicas de \textbf{pessoas} e comunidades, o \textbf{que} fazemos nos negócios afeta o clima e a \textbf{todos} nós.
Outros aspectos envolvem as desigualdades sociais e a maior participação de \textbf{pessoas} \textbf{que} não \textbf{são} brancas e não falam inglês.
Há medo e \textbf{é} preciso coragem, \textbf{é} difícil e então há uma maneira de ter conversas e experiências incríveis e \textbf{reconhecer} \textbf{que} \textbf{cada} um de nós pode estar no centro.
Acho \textbf{que} o COIL vai mudar nos próximos cinco anos se mais \textbf{pessoas} fizerem projetos no ensino fundamental e médio: trabalhar por meio das culturas \textbf{cada} vez mais cedo.
Um dos objetivos \textbf{que} estou buscando \textbf{é} obter informações na comunidade COIL – Centro Paula Souza, SUNY, Florida International University, DePaul, enfim, onde quer \textbf{que} possamos descobrir como fazer crowdsourcing e \textbf{compartilhar} essas informações com \textbf{toda} a comunidade, para \textbf{que} não fiquem isoladas em silos.
Quero \textbf{que} possamos usar big data e extrair informações para os projetos COIL. \textbf{Que} possamos usar hashtags, wikis ou \textbf{formas} de \textbf{compartilhar} informação por meio das comunidades e \textbf{aprender} com isso.
E isso \textbf{é} algo \textbf{que} eu, como designer instrucional, \textbf{sei} \textbf{que} \textbf{é} \textbf{possível}.
Cronologicamente falando, quais foram os avanços mais \textbf{significativos} dos projetos COIL?
No \textbf{início} dos anos 2000, usávamos o Skype e estávamos apenas tentando descobrir como fazer essa \textbf{coisa} de videoconferência e quais ferramentas on-line usar.
Nessa primeira onda, a pergunta \textbf{era}: como fazer?
Depois, entre 2017, 2018 ou mesmo 2019, a pergunta passou a \textbf{ser} o \textbf{que} vamos fazer?
Quando veio a pandemia, as \textbf{pessoas} \textbf{começaram} a se interessar pelo COIL como uma alternativa para \textbf{expandir} as possibilidades de internacionalização para mais estudantes e incluir as \textbf{pessoas} com menos recursos financeiros, como uma comunidade de prática \textbf{que} poderia \textbf{ser} incluída na sala de aula em vez de uma atividade extracurricular.
Viajar, apenas 1 \% consegue.
A \textbf{terceira} onda, de agora em diante, envolve a \textbf{questão}: por \textbf{que} fazemos COIL?
Como abordamos e incluímos \textbf{diferenças} na nossa comunidade, \textbf{aprendendo} sobre os outros e como superamos essas \textbf{diferenças}?
Com a pandemia e o assassinato de George Floyd, \textbf{foi} preciso repensar, reposicionar \textbf{todo} o \textbf{programa} internacional.
O \textbf{que} os \textbf{programas} internacionais e de educação global fornecem de melhor? \textbf{São} conversas sobre cultura e aprendizado experiencial.
Portanto, \textbf{é} preciso \textbf{consolidar} a confiança e buscar o"timing"adequado com base na comunicação.
Isso \textbf{é} o \textbf{que} eu amo tanto no COIL: \textbf{é} uma oportunidade entusiasmante, mesmo quando tantas \textbf{coisas} dão errado.

% matched lemmas: ambiental, aprender, ação, cada, coisa, começar, compartilhar, consolidar, diferença, expandir, forma, início, parte, permitir, pessoa, possível, programa, que, questão, reconhecer, saber, ser, significativo, terceiro, todo
\end{textsample}
